% ============================================================
%  PRACA INŻYNIERSKA
%  Aplikacja do przeprowadzania testów psychologicznych
% ============================================================

\documentclass[12pt, a4paper, oneside]{report}

% --- Kodowanie i język ---
\usepackage[utf8]{inputenc}
\usepackage[T1]{fontenc}
\usepackage[polish]{babel}

% --- Czcionki ---
\usepackage{lmodern}

% --- Geometria strony ---
\usepackage[
  left=3.5cm,
  right=2.5cm,
  top=2.5cm,
  bottom=2.5cm
]{geometry}

% --- Grafika i kolory ---
\usepackage{graphicx}
\usepackage{xcolor}
\usepackage{float}

% --- Tabele ---
\usepackage{booktabs}
\usepackage{array}
\usepackage{longtable}

% --- Odwołania i linki ---
\usepackage[
  colorlinks=true,
  linkcolor=black,
  citecolor=black,
  urlcolor=blue
]{hyperref}

% --- Listingi kodu ---
\usepackage{listings}
\lstset{
  basicstyle=\ttfamily\small,
  breaklines=true,
  frame=single,
  numbers=left,
  numberstyle=\tiny\color{gray},
  keywordstyle=\color{blue},
  commentstyle=\color{green!50!black},
  stringstyle=\color{red!70!black},
  showstringspaces=false,
  tabsize=2
}

% --- Bibliografia ---
\usepackage[
  backend=biber,
  style=numeric,
  sorting=nyt
]{biblatex}
% \addbibresource{bibliografia.bib}  % odkomentuj gdy dodasz plik .bib

% --- Interlinia ---
\usepackage{setspace}
\onehalfspacing

% --- Nagłówki rozdziałów ---
\usepackage{titlesec}
\titleformat{\chapter}[display]
  {\normalfont\Large\bfseries}
  {\chaptertitlename\ \thechapter}{16pt}{\LARGE}

% --- Pakiety dodatkowe ---
\usepackage{enumitem}
\usepackage{caption}
\usepackage{subcaption}
\usepackage{amsmath}
\usepackage{pdfpages}
\usepackage{csquotes}
% Fix for microtype footnote patch warning (conflict with babel/polish + report class)
\usepackage[patch=none]{microtype}
\microtypesetup{nopatch=footnote}


% ============================================================
%  POCZĄTEK DOKUMENTU
% ============================================================
\begin{document}

% ============================================================
%  STRONA TYTUŁOWA
% ============================================================
\begin{titlepage}
  \centering

  \vspace*{1cm}

  {\large\textbf{Uniwersytet Kazimierza Wielkiego w Bydgoszczy}}\\[0.3cm]
  {\large Wydział Informatyki}\\[2cm]

  \rule{\textwidth}{0.4pt}\\[0.6cm]

  {\LARGE\bfseries
    Aplikacja do przeprowadzania testów psychologicznych
  }\\[0.4cm]

  \rule{\textwidth}{0.4pt}\\[2cm]

  {\large Praca inżynierska}\\[2cm]

  \vfill

  {\large Bydgoszcz, 2027}

\end{titlepage}


% ============================================================
%  PODZIĘKOWANIA (opcjonalnie)
% ============================================================
% \chapter*{Podziękowania}
% \addcontentsline{toc}{chapter}{Podziękowania}
%
% TODO: Opcjonalne podziękowania
%
% \newpage


% ============================================================
%  STRESZCZENIE
% ============================================================
%\chapter*{Streszczenie}
%\addcontentsline{toc}{chapter}{Streszczenie}

% TODO: Streszczenie pracy (ok. 0,5 strony)
% Powinno zawierać: cel pracy, zastosowane metody, główne wyniki, wnioski.

%\vspace{1cm}

%\noindent\textbf{Słowa kluczowe:} platforma desktopowa, testy psychologiczne, Electron, Firebase, synchronizacja w chmurze, autoryzacja, RBAC, diagnostyka wspomagana komputerowo

%\newpage


% ============================================================
%  ABSTRACT (wersja angielska)
% ============================================================
%\chapter*{Abstract}
%\addcontentsline{toc}{chapter}{Abstract}

% TODO: Streszczenie w języku angielskim

%\vspace{1cm}

%\noindent\textbf{Keywords:} desktop platform, psychological tests, Electron, Firebase, cloud synchronization, authorization, RBAC, computer-aided diagnostics

%\newpage


% ============================================================
%  SPIS TREŚCI
% ============================================================
\tableofcontents
\newpage


% ============================================================
%  SPIS RYSUNKÓW I TABEL (opcjonalnie)
% ============================================================
%\listoffigures
%\addcontentsline{toc}{chapter}{Spis rysunków}
%\newpage

%\listoftables
%\addcontentsline{toc}{chapter}{Spis tabel}
%\newpage


% %%%%%%%%%%%%%%%%%%%%%%%%%%%%%%%%%%%%%%%%%%%%%%%%%%%%%%%%%%%%
%  ROZDZIAŁ 1 — WSTĘP
% %%%%%%%%%%%%%%%%%%%%%%%%%%%%%%%%%%%%%%%%%%%%%%%%%%%%%%%%%%%%
\chapter{Wstęp}
\label{chap:wstep}

% TODO: Ogólne wprowadzenie w tematykę pracy — kontekst, motywacja.
% Można wspomnieć o rosnącym zapotrzebowaniu na cyfrowe narzędzia
% diagnostyczne w psychologii, potrzebie standaryzacji i bezpieczeństwa danych.


% ============================================================
\section{Cel i zakres pracy}
\label{sec:cel}

Celem niniejszej pracy inżynierskiej jest zaprojektowanie i~zaimplementowanie modularnej platformy desktopowej typu \textit{launcher}, umożliwiającej pobieranie i~uruchamianie cyfrowych testów psychologicznych. Platforma wykorzystuje architekturę klient-serwer, w~której backend stanowi usługa Firebase, a~aplikacja kliencka została zbudowana w~technologii Electron.

Zakres pracy obejmuje:
\begin{itemize}[noitemsep]
  \item Analizę istniejących rozwiązań w~dziedzinie komputerowo wspomaganej diagnostyki psychologicznej.
  \item Zaprojektowanie architektury systemu, w~tym bezpiecznej komunikacji międzyprocesowej (IPC), kontroli dostępu opartej na rolach (RBAC) oraz mechanizmów zapewnienia integralności danych lokalnych.
  \item Implementację aplikacji klienckiej obsługującej rejestrację i~logowanie użytkowników, izolowane uruchamianie testów oraz warunkową synchronizację wyników z~bazą Firestore.
  \item Zapewnienie trybu offline z~lokalnym przechowywaniem danych dla niezweryfikowanych użytkowników oraz pełnej synchronizacji i~archiwizacji w~chmurze (NoSQL) po zatwierdzeniu przez administratora.
  \item Weryfikację poprawności działania systemu względem postawionych wymagań.
\end{itemize}


% ============================================================
\section{Układ pracy}
\label{sec:uklad}

Praca składa się z~następujących rozdziałów:

\begin{description}[style=nextline]
  \item[Rozdział 1 — Wstęp]
    Wprowadzenie w~tematykę, cel i~zakres pracy oraz opis jej struktury.

  \item[Rozdział 2 — Analiza istniejących rozwiązań]
    Przegląd dostępnych platform i~narzędzi do przeprowadzania testów psychologicznych, porównanie ich cech oraz uzasadnienie wyboru własnego rozwiązania.

  \item[Rozdział 3 — Podstawy teoretyczne]
    Omówienie diagnostyki wspomaganej komputerowo, precyzji pomiaru czasu reakcji w~środowiskach opartych na Chromium oraz metod zabezpieczania danych medycznych w~chmurze.

  \item[Rozdział 4 — Projektowanie systemu]
    Opis architektury systemu, komunikacji IPC, modelu RBAC, struktury bazy danych oraz mechanizmów zapewnienia integralności danych.

  \item[Rozdział 5 — Implementacja]
    Szczegóły implementacyjne poszczególnych modułów aplikacji, zastosowane technologie i~biblioteki.

  \item[Rozdział 6 — Weryfikacja i~testowanie]
    Sprawdzenie, czy zrealizowany system spełnia wymagania określone w~temacie pracy.

  \item[Rozdział 7 — Podsumowanie i~wnioski]
    Wnioski z~przeprowadzonych prac, ocena stopnia realizacji celów oraz możliwości dalszego rozwoju platformy.
\end{description}


% %%%%%%%%%%%%%%%%%%%%%%%%%%%%%%%%%%%%%%%%%%%%%%%%%%%%%%%%%%%%
%  ROZDZIAŁ 2 — ANALIZA ISTNIEJĄCYCH ROZWIĄZAŃ
% %%%%%%%%%%%%%%%%%%%%%%%%%%%%%%%%%%%%%%%%%%%%%%%%%%%%%%%%%%%%
\chapter{Analiza istniejących rozwiązań}
\label{chap:analiza}


% ============================================================
\section{Przegląd platform do testów psychologicznych}
\label{sec:przeglad}

% TODO: Opisać istniejące platformy, np.:
% - PsyToolkit
% - Inquisit
\begin{enumerate}
  \item PsychoPy\slash Pavlovia
  \item PsychologyToday
\end{enumerate}
 
% - E-Prime
% - Gorilla Experiment Builder
% - PEBL (Psychology Experiment Building Language)
%
% Dla każdej platformy warto opisać:
% - Główne funkcjonalności
% - Model licencyjny (open-source / komercyjny)
% - Wspierane środowiska (web / desktop)
% - Możliwości offline
% - Sposób przechowywania i synchronizacji danych
% - Precyzja pomiaru czasu reakcji


% ============================================================
\section{Porównanie funkcjonalności}
\label{sec:porownanie}

% TODO: Tabela porównawcza istniejących rozwiązań
%
 \begin{table}[H]
   \centering
   \caption{Porównanie istniejących platform do testów psychologicznych}
   \label{tab:porownanie}
   \begin{tabular}{lcccc}
     \toprule
    \textbf{Cecha} & \textbf{Pavlovia} & \textbf{PsychologyToday} & \textbf{Nous} \\
     \midrule
     Tryb offline       &N &N &T  \\
     Open-source        &N &N &T  \\
     Synchronizacja     &N &N &T  \\
     Desktop            &N &N &T  \\
     Precyzja czasowa   &N &N &T  \\
     Prostota użycia    &T &T &T  \\
     Sytem Tagów        &N &N &T   \\
     \bottomrule
   \end{tabular}
 \end{table}


% ============================================================
\section{Uzasadnienie wybranego rozwiązania}
\label{sec:uzasadnienie}

% TODO: Wyjaśnić, dlaczego zdecydowano się na stworzenie własnej platformy
% zamiast wykorzystania istniejącej. Kluczowe argumenty mogą obejmować:
% - Potrzeba kontroli dostępu (RBAC) z zatwierdzaniem przez administratora
% - Wymagania dotyczące bezpieczeństwa danych medycznych
 \begin{itemize}
     \item Możliwość przeprowadzania testów offline
     \item Modularność i możliwość rozbudowy o nowe testy
     \item Specyficzne wymagania uczelni/klienta
     \item Większa kontrola nad wszelkimi funkcjonalnościami
     
 \end{itemize}


% %%%%%%%%%%%%%%%%%%%%%%%%%%%%%%%%%%%%%%%%%%%%%%%%%%%%%%%%%%%%
%  ROZDZIAŁ 3 — PODSTAWY TEORETYCZNE
% %%%%%%%%%%%%%%%%%%%%%%%%%%%%%%%%%%%%%%%%%%%%%%%%%%%%%%%%%%%%
\chapter{Podstawy teoretyczne}
\label{chap:teoria}


% ============================================================
\section{Diagnostyka wspomagana komputerowo}a
\label{sec:diagnostyka}

% TODO: Omówienie roli komputerów w diagnostyce psychologicznej.
% Historia, korzyści, ograniczenia. Referencje do literatury naukowej.


% ============================================================
\section{Precyzja pomiaru czasu reakcji w środowiskach Chromium}
\label{sec:precyzja}

% TODO: Analiza dokładności temporalnej w przeglądarkach/Electron.
% - requestAnimationFrame vs setTimeout/setInterval
% - performance.now() i High Resolution Time API
% - Wpływ garbage collection, renderowania, V-Sync
% - Porównanie z natywnymi rozwiązaniami (PsychoPy, E-Prime)
% - Strategie minimalizacji niedokładności


% ============================================================
\section{Metody zabezpieczania danych medycznych w chmurze}
\label{sec:bezpieczenstwo}

% TODO: Omówienie aspektów bezpieczeństwa:
% - Szyfrowanie danych w spoczynku i w tranzycie
% - RODO / GDPR w kontekście danych medycznych
% - Firebase Security Rules
% - Modele autoryzacji (RBAC, ABAC)
% - Integralność danych — sumy kontrolne, podpisy cyfrowe


% %%%%%%%%%%%%%%%%%%%%%%%%%%%%%%%%%%%%%%%%%%%%%%%%%%%%%%%%%%%%
%  ROZDZIAŁ 4 — PROJEKTOWANIE SYSTEMU
% %%%%%%%%%%%%%%%%%%%%%%%%%%%%%%%%%%%%%%%%%%%%%%%%%%%%%%%%%%%%
\chapter{Projektowanie systemu}
\label{chap:projektowanie}


% ============================================================
\section{Wymagania funkcjonalne i niefunkcjonalne}
\label{sec:wymagania}

\subsection{Wymagania funkcjonalne}

% TODO: Lista wymagań funkcjonalnych, np.:
% - WF1: System umożliwia rejestrację nowych użytkowników.
% - WF2: Administrator zatwierdza lub odrzuca rejestracje.
% - WF3: Użytkownik może pobierać i uruchamiać testy psychologiczne.
% - WF4: Wyniki testów są zapisywane lokalnie.
% - WF5: Po zatwierdzeniu użytkownika wyniki są synchronizowane z chmurą.
% - ...

\subsection{Wymagania niefunkcjonalne}

% TODO: np.:
% - WN1: Aplikacja działa w trybie offline.
% - WN2: Czas odpowiedzi interfejsu < 200 ms.
% - WN3: Dane są szyfrowane.
% - ...


% ============================================================
\section{Architektura systemu}
\label{sec:architektura}

% TODO: Opis ogólnej architektury klient-serwer.
% Diagram architektury (rysunek).
% Warstwa kliencka (Electron), warstwa backendowa (Firebase),
% moduły testowe (web / PsychoPy).

% \begin{figure}[H]
%   \centering
%   % \includegraphics[width=0.9\textwidth]{img/architektura.png}
%   \caption{Architektura systemu}
%   \label{fig:architektura}
% \end{figure}


% ============================================================
\section{Komunikacja międzyprocesowa (IPC)}
\label{sec:ipc}

% TODO: Opis mechanizmu IPC w Electron:
% - Main process vs Renderer process
% - contextBridge i preload scripts
% - Kanały komunikacji (channels)
% - Bezpieczeństwo IPC — walidacja, sanityzacja


% ============================================================
\section{Model kontroli dostępu (RBAC)}
\label{sec:rbac}

% TODO: Opis modelu ról:
% - Administrator
% - Użytkownik zatwierdzony
% - Użytkownik niezatwierdzony
% Diagram stanów konta użytkownika.
% Firebase Custom Claims.


% ============================================================
\section{Struktura bazy danych}
\label{sec:baza}

% TODO: Model danych w Firestore:
% - Kolekcje i dokumenty
% - Schemat danych użytkowników
% - Schemat wyników testów
% - Reguły bezpieczeństwa Firestore

% \begin{figure}[H]
%   \centering
%   % \includegraphics[width=0.8\textwidth]{img/schemat_bazy.png}
%   \caption{Schemat bazy danych Firestore}
%   \label{fig:baza}
% \end{figure}


% ============================================================
\section{Mechanizmy integralności danych lokalnych}
\label{sec:integralnosc}

% TODO: Opis sposobu zabezpieczenia danych przechowywanych lokalnie:
% - Sumy kontrolne (checksums)
% - Podpisy cyfrowe
% - Walidacja przy synchronizacji


% ============================================================
\section{Projekt interfejsu użytkownika}
\label{sec:ui}

% TODO: Makiety (wireframes) lub zrzuty ekranu projektowanego UI.
% Opis głównych widoków:
% - Ekran logowania/rejestracji
% - Panel główny (launcher)
% - Lista testów
% - Panel administracyjny
% - Ekran wyników


% %%%%%%%%%%%%%%%%%%%%%%%%%%%%%%%%%%%%%%%%%%%%%%%%%%%%%%%%%%%%
%  ROZDZIAŁ 5 — IMPLEMENTACJA
% %%%%%%%%%%%%%%%%%%%%%%%%%%%%%%%%%%%%%%%%%%%%%%%%%%%%%%%%%%%%
\chapter{Implementacja}
\label{chap:implementacja}


% ============================================================
\section{Zastosowane technologie}
\label{sec:technologie}

% TODO: Opis stosu technologicznego:
% - Electron (wersja, uzasadnienie wyboru)
% - JavaScript / HTML / CSS
% - Firebase (Authentication, Firestore, Cloud Functions)
% - PsychoPy (dla modułów testowych w Pythonie)
% - Node.js
% - Narzędzia budowania (electron-builder / electron-forge)


% ============================================================
\section{Struktura projektu}
\label{sec:struktura}

% TODO: Opis organizacji kodu źródłowego.
% Drzewo katalogów, podział na moduły.

% \begin{figure}[H]
%   \centering
%   % \includegraphics[width=0.6\textwidth]{img/struktura_projektu.png}
%   \caption{Struktura katalogów projektu}
%   \label{fig:struktura}
% \end{figure}


% ============================================================
\section{Moduł autoryzacji i rejestracji}
\label{sec:impl-auth}

% TODO: Opis implementacji:
% - Rejestracja użytkownika
% - Logowanie (Firebase Authentication)
% - Progresywna autoryzacja — model zatwierdzania przez administratora
% - Obsługa stanów konta (oczekujący, zatwierdzony, zablokowany)


% ============================================================
\section{Moduł dystrybucji testów}
\label{sec:impl-dystrybucja}

% TODO: Opis implementacji launchera:
% - Pobieranie testów z repozytorium / serwera
% - Aktualizacja testów
% - Izolowane uruchamianie testów
% - Zarządzanie zależnościami testów


% ============================================================
\section{Moduł wykonywania testów}
\label{sec:impl-testy}

% TODO: Opis uruchamiania testów:
% - Testy webowe (JS/HTML)
% - Testy PsychoPy (Python)
% - Komunikacja test ↔ launcher (IPC)
% - Zbieranie i zapisywanie wyników


% ============================================================
\section{Moduł synchronizacji z chmurą}
\label{sec:impl-sync}

% TODO: Opis implementacji synchronizacji:
% - Warunkowa synchronizacja (tylko dla zatwierdzonych użytkowników)
% - Tryb offline — lokalne przechowywanie wyników
% - Resolve konfliktów
% - Archiwizacja danych w Firestore (NoSQL)


% ============================================================
\section{Moduł administracyjny}
\label{sec:impl-admin}

% TODO: Opis panelu administratora:
% - Zarządzanie użytkownikami (zatwierdzanie, blokowanie)
% - Przeglądanie wyników
% - Zarządzanie testami


% %%%%%%%%%%%%%%%%%%%%%%%%%%%%%%%%%%%%%%%%%%%%%%%%%%%%%%%%%%%%
%  ROZDZIAŁ 6 — WERYFIKACJA I TESTOWANIE
% %%%%%%%%%%%%%%%%%%%%%%%%%%%%%%%%%%%%%%%%%%%%%%%%%%%%%%%%%%%%
\chapter{Weryfikacja i testowanie}
\label{chap:weryfikacja}


% ============================================================
\section{Metodyka testowania}
\label{sec:metodyka}

% TODO: Opis podejścia do testowania:
% - Testy jednostkowe
% - Testy integracyjne
% - Testy manualne
% - Środowisko testowe


% ============================================================
\section{Weryfikacja wymagań funkcjonalnych}
\label{sec:wer-funk}

% TODO: Tabela z wymaganiami funkcjonalnymi i ich statusem realizacji.
%
% \begin{table}[H]
%   \centering
%   \caption{Weryfikacja wymagań funkcjonalnych}
%   \label{tab:wer-funk}
%   \begin{tabular}{llc}
%     \toprule
%     \textbf{ID} & \textbf{Wymaganie} & \textbf{Spełnione} \\
%     \midrule
%     WF1 & Rejestracja użytkowników & Tak \\
%     WF2 & Zatwierdzanie przez administratora & Tak \\
%     WF3 & Pobieranie i uruchamianie testów & Tak \\
%     % ...
%     \bottomrule
%   \end{tabular}
% \end{table}


% ============================================================
\section{Weryfikacja wymagań niefunkcjonalnych}
\label{sec:wer-niefunk}

% TODO: Opis testów wydajności, bezpieczeństwa, użyteczności.


% ============================================================
\section{Zgodność z tematem pracy}
\label{sec:zgodnosc}

% TODO: Wykazanie, że zrealizowany system spełnia wymagania
% określone w karcie pracy dyplomowej:
% 1. Platforma desktopowa — zrealizowana w Electron ✓
% 2. Dystrybucja testów — system launchera ✓
% 3. Przeprowadzanie testów psychologicznych — moduły testowe ✓
% 4. Progresywna autoryzacja — RBAC + zatwierdzanie ✓
% 5. Synchronizacja wyników w chmurze — Firestore ✓


% ============================================================
\section{Znalezione problemy i ich rozwiązania}
\label{sec:problemy}

% TODO: Opis napotkanych problemów technicznych i sposobów ich rozwiązania.
% np. problemy z precyzją czasową, Edge-case'y synchronizacji, itp.


% %%%%%%%%%%%%%%%%%%%%%%%%%%%%%%%%%%%%%%%%%%%%%%%%%%%%%%%%%%%%
%  ROZDZIAŁ 7 — PODSUMOWANIE I WNIOSKI
% %%%%%%%%%%%%%%%%%%%%%%%%%%%%%%%%%%%%%%%%%%%%%%%%%%%%%%%%%%%%
\chapter{Podsumowanie i wnioski}
\label{chap:podsumowanie}


% ============================================================
\section{Wnioski}
\label{sec:wnioski}

% TODO: Podsumowanie osiągniętych wyników:
% - Co udało się zrealizować?
% - Jakie cele zostały spełnione?
% - Jakie kompromisy podjęto?


% ============================================================
\section{Ocena realizacji celów}
\label{sec:ocena}

% TODO: Odniesienie do celów z rozdziału 1.
% Stopień realizacji każdego z celów.


% ============================================================
\section{Dalsze możliwości rozwoju}
\label{sec:rozwoj}

% TODO: Propozycje dalszego rozwoju systemu, np.:
% - Rozszerzenie o nowe moduły testowe
% - Wersja mobilna (React Native / Flutter)
% - Zaawansowana analityka wyników (ML / statystyka)
% - Integracja z zewnętrznymi systemami klinicznymi
% - Certyfikacja jako wyrób medyczny
% - Wielojęzyczność interfejsu
% - Szyfrowanie end-to-end (E2EE)
% - Eksport danych do formatów klinicznych (HL7 FHIR)



% ============================================================
\end{document}
